% =================================================================================================
% roteiro.tex — Roteiro de fala da defesa (15 min)
% Acompanha apresentacao.tex (42 slides, com roteiro no início de cada seção).
%
% Compilar: pdflatex roteiro.tex  (2x, por causa da longtable)
% =================================================================================================

\documentclass[11pt,a4paper]{article}

\usepackage[utf8]{inputenc}
\usepackage[T1]{fontenc}
\usepackage[brazilian]{babel}
\usepackage{lmodern}
\usepackage[margin=2cm]{geometry}
\usepackage{array}
\usepackage{longtable}
\usepackage{booktabs}
\usepackage[table]{xcolor}

\definecolor{acento}{RGB}{0,90,156}

% Linha de seção: título em destaque e tempo total da seção
\newcommand{\secao}[2]{%
  \rowcolor{acento!12}
  \multicolumn{3}{@{}p{\dimexpr 16.1cm+4\tabcolsep\relax}@{}}{\rule{0pt}{2.6ex}\textbf{\color{acento}#1} \hfill \textbf{#2}} \\[2pt]
}
% Tempo do slide + tempo acumulado ao final dele
\newcommand{\tempo}[2]{\textbf{#1}\newline{\scriptsize\color{gray}até #2}}
\newcommand{\cad}[1]{\texttt{#1}}

\setlength{\parindent}{0pt}
\renewcommand{\arraystretch}{1.35}

\begin{document}

\begin{center}
  {\Large\bfseries Roteiro de fala — Defesa do TCC}\\[4pt]
  Autosim: DSL para o Ensino de Autômatos Finitos\\[2pt]
  {\small Tempo total: \textbf{15 minutos} \quad$\cdot$\quad 42 slides}
\end{center}

\vspace{4pt}
{\small
\textbf{Como usar:} a coluna \emph{Tempo} traz o tempo do slide e, abaixo, o tempo acumulado ao
final dele. Use o acumulado para conferir o ritmo durante a fala. As falas foram dimensionadas a
cerca de 140 palavras por minuto. Nos slides de \emph{Roteiro} do início de cada seção, diga só a
frase de transição e avance.
}

\vspace{6pt}

\begin{longtable}{@{}>{\raggedright\arraybackslash}p{2.9cm} >{\raggedright\arraybackslash}p{1.3cm} >{\raggedright\arraybackslash}p{11.9cm}@{}}
  \toprule
  \textbf{Slide} & \textbf{Tempo} & \textbf{Fala} \\
  \midrule
  \endfirsthead
  \toprule
  \textbf{Slide} & \textbf{Tempo} & \textbf{Fala} \\
  \midrule
  \endhead
  \bottomrule
  \endlastfoot

  % ---------------------------------------------------------------------------------------------
  \secao{Abertura}{0:30}
  1 — Capa & \tempo{0:15}{0:15} &
  Boa noite a todos. Nós somos João Thomaz Vieira e Klaus Siegfried Beckmann e vamos apresentar
  o nosso trabalho de conclusão de curso: a Autosim, uma ferramenta baseada em linguagem de
  domínio específico para o ensino de autômatos finitos, orientada pelo professor Diógenes Furlan. \\
  2 — Roteiro & \tempo{0:15}{0:30} &
  A apresentação segue a estrutura da monografia: introdução, a fundamentação sobre autômatos,
  compiladores e Rust, os trabalhos relacionados, o projeto da linguagem, a implementação, os
  resultados e a conclusão. \\

  % ---------------------------------------------------------------------------------------------
  \secao{1. Introdução}{2:20}
  3 — Roteiro & \tempo{0:05}{0:35} &
  Começando pela introdução. \\
  4 — Contexto e motivação & \tempo{0:40}{1:15} &
  Autômatos finitos e linguagens formais são um dos pilares da Ciência da Computação: é onde o
  aluno conhece máquinas abstratas, os limites da computação e a base dos compiladores. Mesmo
  assim, é uma disciplina conhecida por ser difícil. Morazán e Antunez apontam o porquê: na
  programação, o aluno escreve, executa e vê o resultado na hora; em autômatos, ele projeta e
  prova a correção no papel, sem poder experimentar nem receber retorno imediato. Isso afasta o
  formalismo da prática. \\
  5 — O problema & \tempo{0:35}{1:50} &
  Existem ferramentas para ajudar, e a mais conhecida é o JFLAP. Mas elas são exclusivamente
  visuais: o autômato é montado arrastando e soltando estados. Com isso, a descrição não é
  texto, o que dificulta versionar e compartilhar; não há validação formal da estrutura; e,
  quando o aluno erra, não recebe um diagnóstico preciso. A lacuna é uma ferramenta que junte o
  rigor da descrição textual, o diagnóstico de um compilador e a interatividade da simulação
  visual. \\
  6 — A proposta: Autosim & \tempo{0:35}{2:25} &
  Nossa proposta é a Autosim, que acopla três partes: uma DSL textual, declarativa e em
  português, para descrever AFDs e AFNs; um compilador pedagógico, com análise léxica,
  sintática e semântica, que aponta os erros com linha e coluna; e uma interface gráfica que
  desenha o diagrama a partir do texto e anima a simulação passo a passo. A contribuição
  central está nesse acoplamento: o aluno escreve o texto e a interface é gerada
  automaticamente, sincronizada com ele. \\
  7 — Objetivos & \tempo{0:25}{2:50} &
  O objetivo geral é integrar descrição textual, compilação pedagógica e simulação visual em um
  único ambiente. Para isso, projetamos a DSL com gramática em EBNF, implementamos o
  compilador, desenvolvemos o simulador, construímos a interface gráfica e validamos tudo com
  autômatos clássicos e cenários de erro. \\

  % ---------------------------------------------------------------------------------------------
  \secao{2. Autômatos Finitos e Linguagens Formais}{1:35}
  8 — Roteiro & \tempo{0:05}{2:55} &
  Agora, a fundamentação, começando pelos autômatos. \\
  9 — Autômatos finitos & \tempo{0:35}{3:30} &
  Formalmente, um AFD é uma quíntupla: um conjunto finito de estados, um alfabeto, uma função de
  transição, um estado inicial e um conjunto de estados de aceitação. A linguagem reconhecida é
  formada pelas cadeias que, partindo do estado inicial, terminam em um estado final. No AFN, a
  transição leva a um conjunto de estados e pode haver transições vazias, as transições
  épsilon. Os dois têm o mesmo poder: reconhecem exatamente as linguagens regulares. \\
  10 — Exemplo prático & \tempo{0:30}{4:00} &
  No diagrama, os círculos são os estados, a seta de entrada marca o inicial, o círculo duplo é
  o final, e as arestas são a função de transição, que também está na tabela ao lado. Esse
  autômato aceita as cadeias que contêm \cad{ab}: aceita \cad{ab} e \cad{aab}, rejeita \cad{ba}
  e a cadeia vazia. Vamos reencontrar esse mesmo exemplo na DSL e na simulação. \\
  11 — Hierarquia de Chomsky & \tempo{0:25}{4:25} &
  Esses autômatos se encaixam na Hierarquia de Chomsky, que classifica as linguagens pela forma
  das regras das gramáticas. No tipo 3 estão as gramáticas regulares, que geram exatamente as
  linguagens reconhecidas pelos autômatos finitos. É nesse nível que o trabalho atua. \\

  % ---------------------------------------------------------------------------------------------
  \secao{3. Compiladores e Linguagens de Domínio Específico}{1:00}
  12 — Roteiro & \tempo{0:05}{4:30} &
  Em seguida, compiladores e linguagens de domínio específico. \\
  13 — Fases de análise & \tempo{0:30}{5:00} &
  Como a Autosim tem um compilador, vale lembrar as fases de análise. A léxica agrupa os
  caracteres em \emph{tokens}, usando expressões regulares. A sintática organiza esses
  \emph{tokens} em uma árvore sintática abstrata, a AST. E a semântica verifica o que a gramática
  sozinha não captura, como referências a algo que não foi declarado. \\
  14 — DSL & \tempo{0:25}{5:25} &
  Uma DSL, segundo Fowler, é uma linguagem de expressividade limitada, voltada a um domínio
  específico: perde generalidade, mas ganha concisão e legibilidade. Ela pode ser interna,
  embutida em outra linguagem, ou externa, com sintaxe própria e compilador dedicado. A Autosim
  é uma DSL externa. \\

  % ---------------------------------------------------------------------------------------------
  \secao{4. Rust como Linguagem de Implementação}{0:30}
  15 — Roteiro & \tempo{0:05}{5:30} &
  Sobre a linguagem de implementação: \\
  16 — Por que Rust? & \tempo{0:25}{5:55} &
  Escolhemos Rust porque ele tem desempenho comparável a C e C++, com segurança de memória
  garantida na compilação e sem coletor de lixo. O sistema de tipos ajuda a representar
  \emph{tokens}, nós da AST e estados sem inconsistências, e o ecossistema, com o Cargo, já
  oferece bibliotecas maduras para \emph{lexing}, \emph{parsing} e interface gráfica. \\

  % ---------------------------------------------------------------------------------------------
  \secao{5. Trabalhos Relacionados}{1:05}
  17 — Roteiro & \tempo{0:05}{6:00} &
  Agora, o que já existe. \\
  18 — Panorama & \tempo{0:30}{6:30} &
  Analisamos cinco trabalhos. O JFLAP é a referência mundial e cobre toda a hierarquia, mas é só
  gráfico. O GAM, da UFLA, é limitado a 10 estados e 3 símbolos. O simulador da UFSC roda na web,
  mas a construção continua puramente gráfica. O AutomatoGraph gera a gramática regular, mas a
  simulação é só parcialmente visual. E o trabalho de Bruzzone, em Rust, tem rigor de
  engenharia, mas não é voltado ao ensino. \\
  19 — Análise comparativa & \tempo{0:30}{7:00} &
  Colocando lado a lado: nenhuma das ferramentas voltadas ao ensino tem descrição textual,
  versionamento ou compilador com diagnóstico; e a única com versionamento não é voltada ao
  ensino nem tem simulação visual. A Autosim é a única que reúne as cinco características. É
  essa lacuna que o trabalho ocupa. \\

  % ---------------------------------------------------------------------------------------------
  \secao{6. Projeto da DSL}{1:25}
  20 — Roteiro & \tempo{0:05}{7:05} &
  Passando para o projeto da linguagem. \\
  21 — A linguagem: um exemplo & \tempo{0:35}{7:40} &
  Este é o mesmo autômato do exemplo anterior, escrito na Autosim. O arquivo \cad{.asl} começa
  declarando o alfabeto. Depois vem o autômato: o tipo, AFD, o nome, os estados, o inicial, os
  finais e as transições, cada uma no formato origem, seta, destino, ``com'' símbolo. No fim, os
  comandos \cad{simular} testam cadeias: \cad{ab} é aceita e \cad{bbb} é rejeitada. A sintaxe é
  em português e próxima da definição formal. \\
  22 — Gramática (EBNF) & \tempo{0:20}{8:00} &
  A linguagem foi formalizada em EBNF. Ela é livre de contexto e pode ser analisada por um
  \emph{parser} descendente preditivo, olhando um único símbolo à frente. São 11 palavras
  reservadas, e nelas maiúsculas e minúsculas não fazem diferença. \\
  23 — Decisões de projeto & \tempo{0:25}{8:25} &
  O alfabeto é declarado uma vez só, o que reforça o conceito teórico e simplifica a
  verificação. O tipo, AFD ou AFN, é explícito, para checar determinismo só no AFD. A
  sintaxe é em português, para o foco ficar nos conceitos, e as palavras-chave não
  diferenciam maiúsculas, para não atrapalhar o aluno. E há suporte a
  transições épsilon com as palavras \cad{epsilon} ou \cad{eps}. \\

  % ---------------------------------------------------------------------------------------------
  \secao{7. Implementação}{2:50}
  24 — Roteiro & \tempo{0:05}{8:30} &
  Agora, a implementação. \\
  25 — Pipeline de compilação & \tempo{0:25}{8:55} &
  O arquivo \cad{.asl} passa pelo analisador léxico, feito com a biblioteca Logos, pelo
  sintático, com Chumsky, e pelo semântico. Se estiver tudo certo, segue para o simulador e
  para a interface, feita com Iced. Se alguma fase encontrar erro, a execução para ali e os
  diagnósticos são exibidos com a biblioteca Ariadne. \\
  26 — Análise léxica e sintática & \tempo{0:25}{9:20} &
  No léxico, o Logos gera um AFD a partir das anotações do tipo \cad{Token}; ou seja, usamos um
  autômato para ler a linguagem de autômatos. Espaços e comentários são descartados, e um
  caractere inválido é reportado com linha e coluna. No sintático, o Chumsky usa \emph{parser
  combinators}: a gramática é escrita compondo funções, e o resultado é a AST. \\
  27 — Análise semântica & \tempo{0:30}{9:50} &
  A semântica tem três fases. Primeiro, a tabela de símbolos, que detecta estados, símbolos ou
  autômatos duplicados. Depois, a verificação de referências: estado não declarado, símbolo
  fora do alfabeto ou simulação de um autômato que não existe. Por último, só para AFD, o
  determinismo: sem épsilon e com uma única transição por estado e símbolo. Os erros são
  acumulados, então o aluno vê todos de uma vez. \\
  28 — Simulador & \tempo{0:35}{10:25} &
  O simulador aplica a função de transição sobre a cadeia e guarda todas as configurações
  intermediárias. Em vez de mostrar a nuvem de estados simultâneos do AFN, ele mostra um caminho
  único: se a cadeia é aceita, o caminho que aceita; se é rejeitada, o caminho que chegou mais
  perto da aceitação. Esse caminho é escolhido por uma busca em largura. E, como o histórico fica
  guardado, o aluno pode avançar e voltar na simulação. \\
  29 — Interface gráfica & \tempo{0:25}{10:50} &
  A interface, feita com Iced, desenha o diagrama automaticamente a partir do texto. A simulação
  pode ser passo a passo ou automática, a 800 milissegundos por passo, destacando o estado
  atual, a transição usada e o que já foi lido. Dá para arrastar e dar zoom no diagrama, e o
  painel fica verde quando a cadeia é aceita e vermelho quando é rejeitada. \\
  30 — Diagnóstico de erros & \tempo{0:25}{11:15} &
  Cada mensagem de erro traz linha e coluna, o trecho do código com realce e uma seta no ponto
  do erro, com a descrição em português. Por exemplo: ``token inválido'', ``AFD não pode ter
  epsilon'' ou ``esperava inicial''. Isso troca o tentar e errar visual por um retorno imediato
  e preciso. \\

  % ---------------------------------------------------------------------------------------------
  \secao{8. Resultados}{2:30}
  31 — Roteiro & \tempo{0:05}{11:20} &
  Vamos aos resultados. \\
  32 — Interface: visão geral & \tempo{0:20}{11:40} &
  Esta é a interface com o autômato do exemplo carregado. No topo, escolhemos a simulação, aqui
  com a cadeia \cad{ab}. O diagrama foi gerado direto do arquivo; embaixo aparecem a entrada,
  com o próximo símbolo destacado, e os botões Voltar, Passo, Rodar, Reset e Centralizar. \\
  33 — Simulação: aceitação & \tempo{0:25}{12:05} &
  Executando a simulação de \cad{ab}: o autômato sai de $q_0$, lê \cad{a} e vai para $q_1$, lê
  \cad{b} e chega em $q_2$, que é final. A cadeia é aceita e o painel fica verde. \\
  34 — Simulação: rejeição & \tempo{0:20}{12:25} &
  Já com \cad{bbb}, o autômato fica em $q_0$, lendo \cad{b} três vezes. A cadeia termina fora
  de um estado final, então é rejeitada e o painel fica vermelho. Embaixo aparece o histórico de
  cada passo. \\
  35 — AFN-$\varepsilon$ e diagramas grandes & \tempo{0:20}{12:45} &
  A ferramenta também trata AFNs com transições épsilon, à esquerda, e autômatos maiores, como
  este AFD de 14 estados, que se navega arrastando e dando zoom. \\
  36 — Testes automatizados & \tempo{0:20}{13:05} &
  Para validar, escrevemos 51 testes automatizados: 12 do léxico, 15 do sintático, 15 do
  semântico e 9 do simulador. Eles cobrem os \emph{tokens}, a construção da AST e todas as
  verificações semânticas, e todos passaram. \\
  37 — Autômatos clássicos & \tempo{0:20}{13:25} &
  Também simulamos autômatos clássicos do Hopcroft e do Sipser: AFD, AFN e AFN com épsilon, com
  cadeias que deveriam ser aceitas e rejeitadas, incluindo a cadeia vazia. Todos os vereditos
  bateram com o esperado. \\
  38 — Cenários de erro & \tempo{0:20}{13:45} &
  E criamos arquivos com erros propositais, como um aluno faria: erros semânticos, AFD com
  épsilon, AFD não determinístico, palavra-chave faltando e caractere inválido. Em todos, a
  mensagem saiu com linha, coluna, trecho realçado e descrição em português. \\

  % ---------------------------------------------------------------------------------------------
  \secao{9. Conclusão}{1:15}
  39 — Roteiro & \tempo{0:05}{13:50} &
  Para concluir. \\
  40 — Contribuições & \tempo{0:30}{14:20} &
  Em resumo, a Autosim aproxima a descrição formal da visualização. O trabalho entrega uma DSL
  em português para AFDs e AFNs, formalizada em EBNF; um compilador pedagógico completo, com
  diagnósticos localizados; um simulador com histórico e navegação nos dois sentidos; e uma
  interface que anima a simulação. Tudo validado com 51 testes e simulações clássicas, todos
  corretos. \\
  41 — Limitações e trabalhos futuros & \tempo{0:30}{14:50} &
  Como limitações: um único alfabeto por arquivo, não há modo de linha de comando, é um arquivo
  por execução, não há editor embutido nem conversão de AFN para AFD ou minimização. Como
  trabalhos futuros, queremos tratar esses pontos, incluir expressões regulares e gramáticas,
  estender a ferramenta para autômatos com pilha e máquinas de Turing e, principalmente, fazer estudos com estudantes para medir o
  impacto no ensino. \\
  42 — Obrigado & \tempo{0:10}{15:00} &
  Obrigado pela atenção. Ficamos à disposição para as perguntas da banca. \\
\end{longtable}

\end{document}
